\documentclass{article}

% ── NeurIPS style ─────────────────────────────────────────
% Download neurips_2025.sty from https://neurips.cc/Conferences/2025/CallForPapers
% and place it next to this .tex file.
% Options: [preprint] shows author names; [final] for camera-ready;
%          [nonatbib] if you want to use your own bib package.
\usepackage[preprint]{neurips_2025}

% ── Packages ──────────────────────────────────────────────
\usepackage[utf8]{inputenc}   % allow utf-8 input
\usepackage[T1]{fontenc}      % use 8-bit T1 fonts
\usepackage{amsmath,amssymb}
\usepackage{booktabs}         % professional tables
\usepackage{graphicx}         % figures
\usepackage{hyperref}         % clickable cross-refs
\usepackage{url}              % simple URL typesetting
\usepackage{microtype}        % microtypography
\usepackage{xcolor}           % colors (used by hyperref)

% ── Bibliography ──────────────────────────────────────────
% natbib is loaded by neurips_2025.sty by default.
% If you pass [nonatbib], uncomment the next line:
% \usepackage[numbers,sort&compress]{natbib}

% ══════════════════════════════════════════════════════════
\title{Opponent-Modulated Value Networks for Exploiting\\
       Suboptimal Play in Leduc Hold'em}

% ── Authors (NeurIPS format) ──────────────────────────────
% Replace with your actual names and affiliations.
\author{%
  Author One \\
  Department of Computer Science \\
  University Name \\
  \texttt{author1@university.edu} \\
  \And
  Author Two \\
  Department of Computer Science \\
  University Name \\
  \texttt{author2@university.edu} \\
  % \And
  % Author Three \\
  % ...
}

\begin{document}
\maketitle

% ══════════════════════════════════════════════════════════
\begin{abstract}
We introduce an opponent-modulated value network for exploiting suboptimal play in Leduc Hold'em. While Counterfactual Regret Minimization (CFR) methods achieve superhuman performance by approximating Nash Equilibrium strategies, they cannot fully capitalize on exploitable opponent tendencies. Our method augments a self-play--trained base value network with a value-modulation head that leverages online opponent statistics to inform decision-making. In round-robin tournaments against a diverse pool of opponents, our approach demonstrates consistent performance gains over CFR-based agents and superior robustness to out-of-distribution opponents compared to standard reinforcement learning baselines.
\end{abstract}

% ══════════════════════════════════════════════════════════
\section{Introduction}

Recent advancement in poker AI has demonstrated superhuman performance across various versions of Texas Hold'em, including Heads-Up Limit~\citep{ref1}, Heads-Up No-Limit~\citep{ref2,ref3,ref6}, and Six-Player No-Limit~\citep{ref4}. These methods predominantly leverage Counterfactual Regret Minimization (CFR) techniques to approximate and play a Nash Equilibrium strategy during test-time. In other words, the policy computed implicitly assumes that opponents play near-optimally, which is valid in matchups against professional poker players.

However, in more casual and realistic settings of gameplay, opponents are not near-optimal. Instead, they exhibit exploitable tendencies, such as over-folding or excessive aggression. When played against suboptimal opponents, CFR-based approaches achieve non-negative expected value (i.e., guaranteed bounded loss), yet they cannot fully capitalize on exploitable patterns. This leads to the central question of this work: To what extent does the equilibrium assumption limit performance against suboptimal, real-world opponents, and how does opponent-aware strategies address this limitation?

We introduced an opponent-modulated value network, utilizing a pre-trained base value network and a value-modulation head for online opponent modeling. During multiple rounds of gameplay against the same opponent, our agent collects opponent statistics, and leverages it via the value-modulation head to inform better decision-making. We evaluated our method against CFR-based agents in a round-robin tournament with a diverse pool of opponents and playing styles, demonstrating that adaptive exploitation yields higher average chip gain per round than equilibrium-seeking strategies.

% ══════════════════════════════════════════════════════════
\section{Background}

\subsection{Information Sets and Definition of Game State}

A core characteristic of imperfect information game is that the full game state isn't observable by players. In poker, the complete state of a game consists of players' private cards, public card (if dealt), pot size, etc. However, players can only observe their own private hand. In other words, the same game state that a player observes (e.g., private card is K and public card is J) can correspond to multiple actual game states (e.g., opponent holding J, Q, or K). This set of game states that a player, given their limited information, cannot distinguish between, is defined as an information set. In principle, a player's policy is a function of the information set they are in, not the actual game state, because the latter is unknown and unobservable.

\subsection{Counterfactual Regret Minimization (CFR)}

CFR is an iterative algorithm for approximating the Nash Equilibrium strategy of an incomplete information game. At every information set of a player, the regret of each available action (under the current iteration's strategy profile) is determined. The accumulated regret over existing iterations is used to construct next iteration's strategy profile, and the average strategy profile over sufficient iterations is theoretically guaranteed to converge to a Nash Equilibrium strategy profile.

CFR cannot be directly applied to poker, because the ${\sim}10^{14}$ information sets is infeasible to compute exactly. Libratus and Pluribus utilize card and action abstraction, grouping similar information sets and actions together to reduce computation, and apply Monte Carlo CFR~\citep{ref5} on the abstracted version of the game. DeepStack and ReBel uses a value function to approximate counterfactual values at a leaf node, limiting the search depth and computation by avoiding to search until a terminal node is reached.

However, CFR assumes all players follow the same iterative procedure when acting, and it outputs a strategy that is optimal only when opponents are also playing Nash Equilibrium strategies. In other words, existing CFR-based methods assume optimal gameplay from opponent, and cannot utilize or exploit opponent behavior that deviates from optimal gameplay.

\subsection{Leduc Hold'em}

In this work, we studied a variant of poker widely used in poker game-theory research, called Leduc Hold'em. Leduc Hold'em consists of a 6-card deck (2 Jacks, 2 Queens, 2 Kings), 2 betting rounds (preflop and flop), 2 players, and 3 possible actions (call, raise, and fold) at each decision point. A detailed walkthrough of Leduc Hold'em rules and gameplay is included in Appendix~\ref{app:rules}.

Leduc Hold'em preserves the core of poker --- imperfect information (both players are dealt private hands that opponent cannot observe) and strategic complexity (behavior like bluffing). Hence, results from this work can be extended to or adapted in other poker variants, such as Texas Hold'em. Furthermore, Leduc Hold'em drastically reduces computation complexity (288 information sets~\citep{ref7}), allowing us to focus on algorithmic innovation without large training overhead. In addition, this means CFR can be used to solve this game exactly, without abstraction, providing a valuable baseline for us to evaluate our result.

% ══════════════════════════════════════════════════════════
\section{Methodology}

Our proposed method consists of 1) a base value network $V_{\text{base}}(s)$ that approximates the expected value of an information set (assuming optimal gameplay onwards), and 2) a value modulation head $\Delta V(s,o)$ that uses collected opponent statistics $o$ to adjust base model's value estimate, optionally through a gating mechanism $g(s,o)$.

\subsection{Base Value Network}

\paragraph{Architecture.}
The base value network's backbone is a Multilayer Perceptron (MLP) with 2 hidden layers of 64 dimensions. We experimented with larger networks, but the increased capacity yielded marginal return. $V_{\text{base}}(s)$ receives a state representation $s\in \mathbb{R}^{15}$ as input, and outputs a scalar that approximates the expected terminal reward of the current player, assuming both players follow a Nash Equilibrium strategy profile from state $s$ onwards. We presented a table illustrating the elements included in our state representation in Appendix~\ref{app:state}.

\paragraph{Self-play.}
We trained the base value network using self-play with a TD(0) loss. During self-play, each agent selects actions via one-step lookahead search. For example, at time step 1 where agent A acts first, it considers the successor states $s_{1,A,\text{call}}$ and $s_{1,A,\text{raise}}$ resulting from each legal action $a_{\text{legal}} \in \mathcal{A}_A$. Each state is evaluated using the value function $V(\cdot)$, and the value estimates are passed into a softmax operator to construct a mixed strategy from which an action is sampled. Formally, agent A selects its action $a$ at its decision point $s_{t,A}$ via
\begin{equation}
a \sim \pi(s_{t,A}) = \mathrm{softmax}\!\bigl(\{V(s_{t,A,a_{\text{legal}}})\}_{a_{\text{legal}}\in \mathcal{A}_A}\bigr)
\end{equation}

Greedy-selection isn't used in test-time since deterministic strategies are easily exploitable in poker.

We presented additional information on alternative training settings (e.g., training against a pool of agents rather than self-play) in the Results section. Empirically, self-play prompts the agent to learn a more robust policy that performs better against out-of-distribution opponents, presumably because self-play introduces a constantly changing opponent, driving the agent toward a mutually near-optimal strategy. In addition, the one-step lookahead formulation we presented is equivalent to learning a Q-network, but it has the additional flexibility to handle continuous action spaces (e.g., arbitrary bet size) since the state after every action can be computed (the state transition function is known). This offers value network formulations better generalization to other poker variants.

\paragraph{Temporal Difference (TD) loss and training.}
We utilized temporal difference loss as the objective function. Specifically, we matched the value estimate of a visited state $s_{t,A,a_{\text{chosen}}}$ to the value of the next visited state $s_{t+2,A,a_{\text{chosen}}}$ ($t{+}2$ since turns alternate and $t{+}1$ is when opponent acts):
\begin{equation}
\mathcal{L}_{\text{TD}(0)} = \mathrm{MSE}\!\bigl(V(s_{t,a_{\text{taken}}}),\; V(s_{t+(0+1)\times 2,\,a_{\text{taken}}})\bigr)
\end{equation}

We also experimented with TD(1) and TD(2), but observed highly unstable training. We found that higher order TD targets frequently reduce to the trajectory's terminal reward (high variance) at later time steps, due to Leduc Hold'em's relatively short time horizon (at most 8 actions before reaching a terminal state). We detailed our training hyperparameters in Appendix~\ref{app:hyper}.

\subsection{Value Modulation Head}

Due to the self-play training dynamic, the base value network implicitly calculates a state's expected value, assuming both players will play the strategy induced by $V_{\text{base}}$ from that state onwards. The value modulation head $\Delta V(s,o)$ adjusts this estimate using observed opponent statistics $o$, producing a composite value estimate:
\begin{equation}
V(s) = V_{\text{base}}(s) + \Delta V(s,o) \cdot g(s,o)
\end{equation}
which approximates state $s$'s expected value, assuming opponent plays a strategy encoded in opponent statistics $o$. In this work, we uses hand-coded features for $o$ rather than training an end-to-end opponent encoder. This potentially simplified the modulation head's task to learning different value adjustments for different playing styles summarized in $o$.

\paragraph{Opponent statistics.}
We introduced the notion of poker sessions, where each session consists of $n=100$ consecutive games played between two agents. As a session progresses, each agent gradually accumulate statistics about opponent, such as preflop fold rate and re-raise rate (details presented in Appendix~\ref{app:oppstats}). We also included a scalar $\text{number of observed rounds}/n$ to represent the reliability of the collected statistics. To mitigate the cold start problem, we utilized a prior (average statistics of agents in training pool) via Beta-Bernoulli smoothing, with prior strength $S=20$.

\paragraph{Architecture and gating mechanism.}
We used a MLP with two 32-dimension hidden layers for the value modulation head. We included an optional gate, conditioned on opponent statistics and the current game state, to control the influence of the modulation. This is motivated by the observation that opponent style has variable impact on different states.

\paragraph{Pool training.}
We trained the value modulation head using similar recipe as the base value network. Instead of self-play, we sample an opponent from the training pool for each session and played our agent against it. The training pool contains 6 hand-crafted rule-based agents representing distinct playing styles, a CFR-based agent, and the base value model. This setup exposes our agent to a wide range of playing styles, allowing it to better learn how different opponent statistics can influence value. In addition, we freeze the base value network's weights during training, because it is not conditioned on opponent statistics, hence it will just overfit the training pool if allowed to update.

% ══════════════════════════════════════════════════════════
\section{Results}

We evaluated our method against baseline models, including CFR, Policy Gradient (REINFORCE), Actor-Critic, and Double Q-Learning (DQN). These models are trained with similar budgets against the same agent pool (In-Distribution) used for training the value modulation head. To evaluate each model, we played it against all agents in the opponent pool (In-Distribution + OOD, details in Appendix~\ref{app:pool}), 10000 games (100 sessions) per opponent. For each opponent (x-axis), we measured the chips/session earned by each model, and subtract from that the chips/session earned by the CFR-based model. This measures the additional value each model extracts by deviating from optimal play.

% ┌──────────────────────────────────────────────────────────┐
% │  FIGURE 1 — main results bar chart                       │
% │  Uncomment and update the filename once you have it.     │
% └──────────────────────────────────────────────────────────┘
% \begin{figure}[t]
%   \centering
%   \includegraphics[width=\linewidth]{figures/main_results.pdf}
%   \caption{Per-opponent chips/session gain over CFR.
%            Our method (bolded purple) stays above zero against
%            almost all opponents.}
%   \label{fig:main}
% \end{figure}

Our method (bolded purple) demonstrated consistent performance gains over CFR-based model, staying above 0 net-gain against almost all opponents. Although baseline methods exhibit better exploitation of rule-based agents (in-distribution), the policy they learned are not as robust, and their performance collapsed against out-of-distribution (OOD) agents.

% ── Table 1: Main results ─────────────────────────────────
\begin{table}[t]
\centering
\caption{Agent performance in chips per 100-hand session on Leduc Hold'em.
\textbf{Rob.}$\;{=}\;\mu - 1.5\sigma$ (higher is better).
Best per column among non-reference agents is \textbf{bolded}.
Self-matchup scores set to 0 by symmetry.}
\label{tab:main}
\setlength{\tabcolsep}{5pt}
\begin{tabular}{l rr rr rr}
\toprule
& \multicolumn{2}{c}{In-Distribution} & \multicolumn{2}{c}{OOD} & \multicolumn{2}{c}{Overall} \\
\cmidrule(lr){2-3}\cmidrule(lr){4-5}\cmidrule(lr){6-7}
\textbf{Agent} & Avg & Rob. & Avg & Rob. & Avg & Rob. \\
\midrule
\textit{CFR (Nash)} & $+25.6$ & $-30.8$ & $+12.3$ & $-16.4$ & $+20.8$ & $-28.4$ \\
\midrule
Value Network & $+44.7$ & $-29.5$ & $-25.9$ & $-25.9$ & $+35.9$ & $-41.8$ \\
\textbf{Modulated Value Net.\ (Ours)} & $+50.1$ & $-31.1$ & $\mathbf{+27.5}$ & $-20.7$ & $\mathbf{+41.9}$ & $-31.0$ \\
REINFORCE & $+53.4$ & $-28.3$ & $+12.8$ & $\mathbf{-2.1}$ & $+38.6$ & $-33.4$ \\
Actor-Critic & $+2.2$ & $-94.3$ & $-27.7$ & $-59.0$ & $-8.7$ & $-90.8$ \\
DQN & $\mathbf{+59.1}$ & $\mathbf{-5.4}$ & $+5.4$ & $-23.1$ & $+39.6$ & $\mathbf{-27.1}$ \\
\bottomrule
\end{tabular}
\end{table}

\subsection{Value Modulation Head Analysis}

We plotted the value modulation head's output (y-axis) at each possible information set (x-axis), for three of the rule-based agents with distinct styles. The modulation head exhibits differentiation across opponent types (e.g., positive modulation against loose-aggressive opponent who raise too often, versus negative modulation against tight-passive). However, it demonstrated limited capability for fine-grained modulation between states, hinting that it mainly learned a useful global adjustment to all states for different opponent types.

% ┌──────────────────────────────────────────────────────────┐
% │  FIGURE 2 — modulation head analysis                     │
% │  Uncomment and update the filename once you have it.     │
% └──────────────────────────────────────────────────────────┘
% \begin{figure}[t]
%   \centering
%   \includegraphics[width=\linewidth]{figures/modulation_analysis.pdf}
%   \caption{Value modulation head output across information sets
%            for three opponent types.}
%   \label{fig:modulation}
% \end{figure}

\subsection{Ablation Study}

To understand the contribution of each component of our method, we conducted an ablation study summarized in Table~\ref{tab:ablation}.

\paragraph{Impact of opponent statistics.}
Compared to base value network, adding modulation using opponent statistics yielded higher average return. To account for the increased capacity brought by the modulation head (a 3 layer MLP), we trained another modulation head that does not use opponent statistics as input (+ State Mod., no opp.\ stats). This produced a modest in-distribution gain but caused a significant drop in OOD performance. This confirmed that value modulation head's pool-training does allow it to learn opponent-value modulation mappings that are generalizable to opponents of unseen playing style.

\paragraph{Impact of freezing the base network.}
Training with an unfrozen base network (+ Opp.\ Stats Mod., unfrozen base) achieves the highest in-distribution average (+62.7), but collapses on OOD opponents ($-17.6$). This supports our hypothesis that an unfrozen base network absorbs pool-specific patterns into its weights, sacrificing generality for in-distribution fit.

\paragraph{Joint training from scratch.}
Training the full composite network end-to-end from scratch performs worst overall (+29.4 avg, $-38.1$ robustness), suggesting that a well-initialized base network is critical --- the modulation head alone cannot compensate for a poorly learned value prior.

% ── Table 2: Ablation ─────────────────────────────────────
\begin{table}[t]
\centering
\caption{Ablation study on Leduc Hold'em (chips per 100-hand session, averaged across seeds).
\textbf{Rob.}~$= \mu - 1.5\sigma$ (higher is better). Best per column in \textbf{bold}.}
\label{tab:ablation}
\setlength{\tabcolsep}{5pt}
\begin{tabular}{l rr rr rr}
\toprule
& \multicolumn{2}{c}{In-Distribution} & \multicolumn{2}{c}{OOD} & \multicolumn{2}{c}{Overall} \\
\cmidrule(lr){2-3}\cmidrule(lr){4-5}\cmidrule(lr){6-7}
Method & Avg & Rob. & Avg & Rob. & Avg & Rob. \\
\midrule
Value Net.\ (self-play) & $+44.7$ & $-29.5$ & $-25.9$ & $-25.9$ & $+35.9$ & $-41.8$ \\
\quad + State Mod.\ (no opp.\ stats) & $+51.8$ & $-24.9$ & $-28.0$ & $-28.0$ & $+41.8$ & $-40.1$ \\
\textbf{\quad + Opp.\ Stats Mod., frozen base}~\textbf{(Ours)} & $+50.1$ & $-31.1$ & $\mathbf{+27.5}$ & $-20.7$ & $+41.9$ & $-31.0$ \\
\quad + Opp.\ Stats Mod., unfrozen base & $\mathbf{+62.7}$ & $\mathbf{-10.1}$ & $-17.6$ & $\mathbf{-17.6}$ & $\mathbf{+52.6}$ & $\mathbf{-26.2}$ \\
\bottomrule
\end{tabular}
\end{table}

% ══════════════════════════════════════════════════════════
% NeurIPS requires references before the appendix.
% The style enforces a 10-page limit for the main body;
% appendices and references do not count toward the limit.
% ══════════════════════════════════════════════════════════
\bibliographystyle{plainnat}
\bibliography{references}

% ══════════════════════════════════════════════════════════
\appendix

\section{Leduc Hold'em Rules}
\label{app:rules}

Leduc Holdem is a simplified version of Poker. As in Poker, at any point in the game, the player can choose from a set of three possible actions [Call, Raise, Fold]. The card deck for Leduc consists of 6 cards [King, King, Queen, Queen, Jack, Jack]. The game is played between exactly two players.

The game has two rounds, pre flop and post flop. Pre flop is before a board card is introduced and post flop is after a board card is introduced. A player can raise 2 chips in the pre-flop round and 4 chips in the post-flop round. In each betting round, a maximum of two raises is allowed per player.

At the beginning of each game, each player adds one chip into the pot. The initial pot has two chips. Then, a private card is dealt to each player. Following this, the pre flop round commences. Each player can choose from the set of possible actions. Either player can choose to call (match the bet), raise (increase the bet) or fold, which would signify the end of the game. After the first round of betting, a board card is introduced, which is visible to both players. Following this, the post flop round commences. The game ends in two ways. First is if, at any point, either player folds, that would terminate the game and the opposing player wins. The other is during a showdown. In this, both players reveal their cards and the player with the matching board card wins. If neither player has a matching card, then the player with the higher priority card wins. The priority order is (King$>$Queen$>$Jack). The entirety of the pot goes to the winning player.

\section{State Representation}
\label{app:state}

\begin{table}[h]
\centering
\begin{tabular}{lll}
\toprule
\textbf{Name} & \textbf{Data Structure} & \textbf{Description} \\
\midrule
Player\_Hands       & String          & Both player's private hands. \\
Board               & String          & Board Card (None until post flop round). \\
Pot                 & List            & 2-element list containing the chips bet by each player. \\
Current\_Round      & Integer         & 0 or 1 to indicate pre flop and post flop respectively. \\
Current\_Player     & String          & Indicates the player whose turn it is. \\
Raises\_this\_Round & Integer         & Counter for total raises in the current round. \\
History             & List of Tuples  & A list of (player, action\_name) tuples. \\
round\_betting\_ended   & Boolean    & Flag: has current round betting ended? \\
last\_action\_was\_raised & Boolean  & Flag: was the most recent action a raise? \\
\bottomrule
\end{tabular}
\caption{Elements of the game state representation.}
\label{tab:state}
\end{table}

\section{Training Hyperparameters}
\label{app:hyper}

\begin{table}[h]
\centering
\begin{tabular}{lrl}
\toprule
\textbf{Hyperparameter} & \textbf{Value} & \textbf{Description} \\
\midrule
Learning rate       & $1 \times 10^{-4}$  & Step size for Adam per batch update \\
Hidden layer size   & 64                   & Neurons per hidden layer \\
Hidden layers       & 2                    & Depth of the value network \\
Optimiser           & Adam                 & $\beta_1{=}0.9,\;\beta_2{=}0.999$ \\
Loss function       & MSE                  & Between predicted and target $V(s)$ \\
Temperature ($\tau$)& 1.0                  & Boltzmann exploration scale \\
Batch size          & 32                   & Episodes per gradient update \\
Eval interval       & 50 episodes          & Checkpoint frequency \\
Eval games          & 100                  & Hands per checkpoint evaluation \\
\bottomrule
\end{tabular}
\caption{Training hyperparameters for the base value network.}
\label{tab:hyper}
\end{table}

\section{Opponent Statistics}
\label{app:oppstats}

The value agent also contains a few derived factors of opponent behaviour:

\begin{description}
\item[Fold Rate:] The ratio of total folds taken to the total number of actions taken. A high fold rate may either indicate a passive opponent or a potential bluff.
\item[Raise Rate:] The ratio of total raises taken to the total number of actions taken. A high raise rate may either indicate an aggressive opponent or a potential bluff.
\item[Fold to Raise Rate:] This is the ratio of instances when the opponent\ldots
\item[Confidence:] \emph{(to be completed)}
\end{description}

\section{Agent Pool}
\label{app:pool}

% TODO: Add agent pool table/description here.

\end{document}
